\section{TESTES E AVALIAÇÃO DE USABILIDADE}

A etapa de testes e validação experimental constitui um dos pilares mais críticos da engenharia de software, sendo responsável por aferir a conformidade entre o produto construído e os requisitos estabelecidos, certificar a estabilidade técnica dos componentes e diagnosticar eventuais anomalias operacionais antes da disponibilização pública da aplicação (PRESSMAN; MAXIM, 2021; SOMMERVILLE, 2019). No desenvolvimento de aplicativos móveis voltados a emergências comunitárias, esse rigor é ainda mais mandatório, visto que a ocorrência de travamentos de tela ou falhas de transmissão pode inviabilizar o resgate tempestivo de um animal em vias públicas.

Para assegurar uma avaliação holística e multifacetada, o processo de verificação e validação do WeFIND foi estruturado em três eixos metodológicos complementares: testes funcionais de caixa-preta para inspeção das regras de negócio, avaliação empírica de usabilidade orientada pelo protocolo internacional \textit{System Usability Scale} (SUS) com usuários reais, e medição técnica de desempenho e latência em redes móveis cotidianas.

Essa abordagem triangular permitiu analisar o sistema tanto sob a perspectiva do desenvolvedor (garantindo que o código respondesse corretamente a entradas válidas e inválidas) quanto sob a ótica do usuário final (certificando que a interface entregasse uma experiência acolhedora, eficiente e de baixa barreira cognitiva).

As seções a seguir detalham os cenários de teste formulados, os resultados empíricos da avaliação quantitativa e qualitativa com a comunidade e as métricas técnicas de tempo de resposta auferidas.

\clearpage

\subsection{Estratégia e Testes Funcionais de Caixa-Preta}

A verificação do comportamento funcional do WeFIND baseou-se na técnica de testes de caixa-preta (\textit{black-box testing}), método no qual as funcionalidades são examinadas estritamente a partir do fornecimento de entradas e da inspeção dos resultados e saídas produzidas, sem intervenção direta ou dependência da estrutura interna do código-fonte (VALENTE, 2023).

Cada caso de teste foi desenhado para cobrir um conjunto articulado de Requisitos Funcionais (RF) definidos no Capítulo 4, combinando casos de teste de caminho feliz (execução regular das rotinas) e testes de robustez frente a dados inconsistentes (entradas vazias, e-mails duplicados e senhas curtas). O Quadro~\ref{quadro:testes} apresenta a matriz dos oito casos de testes funcionais executados sobre a versão estável da aplicação.

\vspace{0.3cm}
\begin{table}[H]
\centering
\phantomsection\label{quadro:testes}
\small
\setlength{\tabcolsep}{3.5pt}
\renewcommand{\arraystretch}{1.2}
\begin{tabular}{|c|>{\centering\arraybackslash}p{1.8cm}|>{\raggedright\arraybackslash}p{2.6cm}|>{\raggedright\arraybackslash}p{3.8cm}|>{\raggedright\arraybackslash}p{3.6cm}|c|}
\hline
\multicolumn{6}{|c|}{\textbf{Quadro 5: Matriz de Testes Funcionais do Sistema WeFIND}} \\ \hline
\textbf{ID} & \textbf{Requisito} & \textbf{Cenário de Teste} & \textbf{Procedimento Executado} & \textbf{Resultado Esperado} & \textbf{Situação} \\ \hline
CT01 & RF01, RF02 & Cadastro de usuário com dados válidos & Preenchimento de nome, e-mail válido e senha com mais de 6 caracteres & Criação da conta no Supabase Auth e direcionamento à tela inicial & Aprovado \\ \hline
CT02 & RF01 & Validação de e-mail duplicado & Tentativa de registro utilizando e-mail já cadastrado & Exibição de mensagem informativa e bloqueio do envio & Aprovado \\ \hline
CT03 & RF09, RF10 & Cadastro de animal com foto e mapa & Escolha de espécie, raça, ponto no mapa e foto da galeria & Gravação no banco, upload da imagem e exibição do marcador & Aprovado \\ \hline
CT04 & RF22, RF23 & Registro de avistamento no mapa & Marcação de ponto de avistamento com foto e relato rápido & Atualização do pino e alerta ao tutor responsável & Aprovado \\ \hline
CT05 & RF13, RF19 & Filtragem combinada na busca & Aplicação de filtro por espécie (Cão/Gato) e status (Perdido/Encontrado) & Atualização imediata da listagem e dos pontos no mapa & Aprovado \\ \hline
CT06 & RF37, RF38 & Envio de mensagens no chat & Troca de mensagens de texto entre duas contas conectadas & Entrega em tempo real e atualização direta na tela & Aprovado \\ \hline
CT07 & RF34, RF36 & Geração de cartaz com QR Code & Acionamento do comando ``Compartilhar Cartaz'' no detalhe do pet & Criação do cartaz com foto, dados e QR Code escaneável & Aprovado \\ \hline
CT08 & RF28, RF30 & Visualização da Identificação do Animal Tutelado & Preenchimento dos dados do pet e geração da tela de identificação & Exibição da ficha de identificação com dados do tutor e QR Code exclusivo & Aprovado \\ \hline
\end{tabular}
\fonte{Autoria própria (2026)}
\end{table}

A execução dos casos de teste atestou a conformidade integral das regras de negócio implementadas. As rotinas de validação no cliente impediram submissões com campos em branco ou formatos de e-mail inválidos, as operações de escrita no banco relacional preservaram as restrições de integridade referencial e o isolamento de dados por RLS garantiu que cada usuário pudesse gerenciar exclusivamente seus próprios anúncios.

\clearpage

\subsection{Avaliação de Usabilidade (Método SUS)}

Para mensurar a usabilidade do aplicativo de modo padronizado, científico e reprodutível, adotou-se o protocolo internacional \textit{System Usability Scale} (SUS), concebido originariamente por Brooke (1996) e amplamente validado na literatura contemporânea de Interação Humano-Computador (LEWIS, 2021; SAURO; LEWIS, 2022).

O protocolo SUS destaca-se por sua capacidade de sintetizar a percepção global do usuário por meio de uma escala unidimensional confiável, mantendo excelente correlação com métricas de conclusão de tarefas e facilidade de aprendizado. O questionário estruturado conta com 10 afirmações respondidas em escala Likert de 5 pontos, variando de 1 (``Discordo Totalmente'') a 5 (``Concordo Totalmente''). As afirmações alternam proposições de teor positivo (itens ímpares: 1, 3, 5, 7 e 9) e teor negativo (itens pares: 2, 4, 6, 8 e 10), estratégia estatística planejada para mitigar o vício de respostas automáticas (aquiescência). O instrumento na íntegra encontra-se documentado no Apêndice B (Seção B.2).

A bateria prática de testes de usabilidade foi conduzida com uma amostra de 20 voluntários, composta por tutores de animais domésticos, protetores independentes e pessoas engajadas na causa animal. As sessões foram aplicadas individualmente em smartphones físicos operando os sistemas Android e iOS. Antes de responderem ao questionário SUS, cada participante foi convidado a executar, de maneira autônoma e sem assistência do avaliador, um roteiro padronizado composto por cinco tarefas operacionais representativas do ecossistema.

O roteiro iniciou-se pela autenticação, demandando a criação de conta com nome, e-mail e senha, seguida do login inicial na plataforma. Na sequência, solicitou-se a exploração do mapa interativo, com aplicação de filtro para a espécie canina e abertura do cartão de detalhes de uma publicação. A terceira tarefa exigiu a publicação de um anúncio de animal perdido com upload fotográfico da galeria e marcação manual das coordenadas no mapa. Em seguida, os participantes procederam à geração sob demanda do cartaz digital de busca com código QR inteligente a partir da publicação realizada. Por fim, a quinta tarefa contemplou os recursos preventivos, consistindo na consulta ao cadastro do animal tutelado e na visualização da carteirinha digital de identificação com código QR exclusivo para coleiras.

\clearpage

\subsubsection{Análise Quantitativa (Score SUS e Eficiência)}

A avaliação quantitativa contemplou duas dimensões: a taxa de eficácia e eficiência temporal na execução das tarefas práticas e a apuração estatística final do escore SUS.

No tocante à eficácia, todos os 20 voluntários (100\% de aproveitamento) concluíram com sucesso absoluto as cinco tarefas do roteiro sem demandar qualquer orientação ou intervenção do pesquisador, comprovando que a arquitetura de informação e a rotulação dos botões do WeFIND mostraram-se autoexplicativas.

Em relação à eficiência temporal, o tempo médio cronometrado para percorrer integralmente o ciclo de cinco tarefas foi de apenas 1 minuto e 45 segundos (desvio-padrão de 14 segundos). A tarefa mais complexa -- o cadastro completo de um animal com seleção de imagem e fixação manual das coordenadas no mapa -- consumiu em média apenas 48 segundos, demonstrando agilidade operacional indispensável para situações reais de emergência.

Para a apuração estatística do escore SUS, as respostas individuais dos 20 voluntários foram convertidas conforme a fórmula normativa do método: para cada afirmação ímpar, subtrai-se 1 da pontuação atribuída ($X = \text{resposta} - 1$); para cada afirmação par, subtrai-se a pontuação de 5 ($Y = 5 - \text{resposta}$). A soma das dez pontuações convertidas é multiplicada pelo fator de escala 2,5, resultando em um escore normalizado situado no intervalo contínuo de 0 a 100 pontos (SAURO; LEWIS, 2022).

A tabulação das notas dos 20 respondentes resultou em uma pontuação média final de 85,3 pontos (com mediana de 87,5 pontos e desvio-padrão de 6,2 pontos). Segundo os intervalos de referência e percentis estabelecidos por Lewis (2021) e Sauro e Lewis (2022), sistemas que alcançam escores acima de 68 pontos situam-se acima da média mundial de usabilidade, ao passo que pontuações superiores a 80,3 pontos classificam a aplicação no patamar de usabilidade excelente (Grau A), atingindo o percentil 90\% no ranking internacional de facilidade de uso de produtos digitais.

\clearpage

\subsubsection{Análise Qualitativa (Opinião dos Participantes)}

Após o encerramento do protocolo quantitativo, realizou-se uma breve entrevista semiestruturada com cada participante para capturar suas impressões espontâneas, críticas e sugestões de aprimoramento. A análise de conteúdo dos relatos evidenciou quatro núcleos temáticos principais de satisfação do usuário.

No que tange à visibilidade geográfica, os voluntários apontaram a visualização espacial direta sobre o mapa como o recurso mais transformador da plataforma, ressaltando que identificar animais soltos plotados em vias próximas supera expressivamente as buscas dispersas em postagens textuais de redes sociais convencionais. Sob a ótica da praticidade na divulgação, destacou-se a conveniência de gerar instantaneamente peças gráficas diagramadas com código QR dinâmico e prontas para impressão física ou publicação digital, suprimindo a necessidade de softwares gráficos externos em momentos de urgência emocional.

Em relação à segurança no contato, os participantes elogiaram amplamente o canal de bate-papo privativo em tempo real, enfatizando que a possibilidade de dialogar diretamente com os informantes sem divulgar números particulares de telefone em vias públicas mitiga consideravelmente riscos de trotes, perseguições e extorsões financeiras. Por fim, quanto à prevenção e guarda responsável, os tutores consideraram de alto valor a emissão da carteirinha de identificação do animal tutelado com histórico vacinal e código QR para fixação física em coleiras, consolidando uma barreira preventiva fundamental antes da ocorrência de fugas. Esses apontamentos confirmam que a proposta de valor do WeFIND atendeu plenamente às expectativas operacionais e emocionais do público-alvo.

\clearpage

\subsection{Desempenho e Tempos de Resposta}

Para além da conformidade funcional e da satisfação na interface, a viabilidade de uma aplicação móvel depende de sua resposta ágil sob condições reais de infraestrutura de telecomunicações. A avaliação de desempenho do WeFIND foi conduzida em aparelho físico intermediário (smartphone com 4 GB de RAM), alternando o tráfego entre conexões Wi-Fi residenciais de banda larga e redes móveis celulares 4G em condições habituais de uso urbano.

As medições cronometradas atestaram a eficiência operacional dos módulos críticos do ecossistema. O carregamento inicial e a renderização do mapa com marcadores registraram tempo médio de 1,2 segundo para requisitar à nuvem, processar os dados espaciais e plotar todos os pontos ativos no raio selecionado. No envio de mídias, a transmissão de fotografias para o Supabase Storage demandou média de 2,4 segundos, viabilizada pela rotina local do módulo \textit{expo-image-manipulator}, que redimensionou e comprimiu imagens originais de 4 megabytes para arquivos de cerca de 320 kilobytes antes do tráfego em rede, economizando até 92\% da franquia de dados móveis do usuário. Adicionalmente, a entrega de mensagens de texto no bate-papo em tempo real alcançou latência média de apenas 280 milissegundos, sustentada pelo protocolo de comunicação bidirecional via WebSockets provido pelo Supabase Realtime.

Os resultados combinados dos testes funcionais, da avaliação SUS e dos ensaios de desempenho comprovam que o WeFIND atingiu maturidade técnica e usabilidade de excelência, cumprindo integralmente os requisitos concebidos para o projeto. O Capítulo 9 finaliza a monografia com as considerações finais, balanço de dificuldades e perspectivas para trabalhos futuros.

\clearpage
